\chapter{WASI Preview 1: WasmEdge Historical Comparison}
\label{app:wasmedge_p1}

This appendix documents the \acs{WASI}~P1 implementation and benchmark results obtained using the \acs{WasmEdge} runtime. This work was completed as a first phase of the experiment before the primary runtime was migrated to SpinKube/Wasmtime (\acs{WASI}~P2). The results here serve as a historical baseline and provide concrete quantitative evidence for the \acs{WASI}~P1 concurrency constraints that motivated the migration.

To reproduce the \acs{WASI}~P1 environment, set \texttt{ENABLE\_WASMEDGE=true} in \texttt{cloud-init.sh}, run \texttt{make label-node-wasmedge} and \texttt{make deploy-wasmedge} in \texttt{thesis-infra-setup}, and apply the optional manifests in \texttt{k8s/01-prime-sieve/optional/}. Source code is in \texttt{wasm/wasmedge/}.

% ── Section 1: Why WasmEdge Was the Original Runtime ─────────────────────────
\section{Why WasmEdge Was the Original Runtime}
\label{app:wasmedge_rationale}

When the experiment was first designed, the research question required a \acs{Wasm} runtime that could serve as a persistent \acs{HTTP} microservice inside a Kubernetes pod---architecturally comparable to a long-running Docker container that owns its own \acs{TCP} listener. This constraint eliminated most runtimes and forced the selection of \acs{WasmEdge}~0.14.1.

The four binding constraints were:

\begin{enumerate}
    \item \textbf{No async Rust HTTP frameworks for \texttt{wasm32-wasip1}.}
          \texttt{hyper\_wasi} and \texttt{tokio\_wasi} are unmaintained \acs{WasmEdge}-specific forks that fail to compile against Rust $\geq$1.80. \texttt{std::net::TcpListener::bind("0.0.0.0:8080")} returns \texttt{Err("Unsupported")} at runtime on \texttt{wasm32-wasip1}: \acs{WASI}~P1 does not expose \acs{POSIX} socket primitives natively. The only viable crate was \texttt{wasmedge\_wasi\_socket~0.5.5}~\cite{wasmedgewasisocket}, a \acs{WasmEdge}-specific extension.

    \item \textbf{TinyGo's \texttt{net/http} server is broken on \texttt{wasip1} since v0.28.}
          The Netdev architecture refactor in TinyGo~0.28 permanently disabled the server-side \texttt{ListenAndServe} code path for the \acs{WASI} target. The only solution was a custom \texttt{server.go} calling \acs{WasmEdge}'s \acs{TCP} socket host functions via Go's \texttt{//go:wasmimport} directive.

    \item \textbf{The \texttt{runwasi} musl-static shim cannot support networking.}
          The \texttt{containerd/runwasi} project provides a pre-built musl-static \acs{WasmEdge} shim, but it bundles its own \acs{WasmEdge} without the proprietary socket extension compiled in. Using the \texttt{runwasi} shim would silently drop socket calls. A custom crun~v1.22 compiled with \texttt{--with-wasmedge} was therefore required.

    \item \textbf{Wasmtime's \acs{WASI}~P1 lacks TinyGo's socket ABI.}
          Wasmtime implements only the standardised \acs{WASI} syscalls; it does not implement \acs{WasmEdge}'s proprietary \texttt{sock\_open}/\texttt{sock\_bind}/\texttt{sock\_listen} extensions. TinyGo's \texttt{//go:wasmimport} directives call these \acs{WasmEdge}-specific functions, making TinyGo's \acs{WASI}~P1 networking Wasmtime-incompatible.
\end{enumerate}

These constraints made \acs{WasmEdge} the only runtime that simultaneously satisfied all experimental requirements in the \acs{WASI}~P1 era. They also directly motivated the migration to SpinKube/\acs{WASI}~P2, where none of these workarounds are needed.

% ── Section 2: WasmEdge Runtime Chain ────────────────────────────────────────
\section{WasmEdge Runtime Chain}
\label{app:wasmedge_chain}

Running \acs{Wasm} pods alongside standard \acs{OCI} pods requires configuring containerd to route \acs{Wasm} workloads through a different \acs{OCI} runtime. The full \acs{WasmEdge} chain is:

\begin{center}
\texttt{kubelet} $\to$ \texttt{containerd} $\to$ \textit{wasmedge shim handler} $\to$ \texttt{crun~v1.22} (\texttt{--with-wasmedge}) $\to$ \texttt{WasmEdge~0.14.1} $\to$ \texttt{.wasm binary}
\end{center}

The key components in this chain are:

\begin{itemize}
    \item \textbf{crun~v1.22:} A lightweight \acs{OCI} container runtime written in C. Unlike \texttt{runc} (written in Go), crun can be compiled with \texttt{--with-wasmedge} to link against \texttt{libwasmedge.so}, enabling it to detect and execute \texttt{.wasm} binaries natively. No pre-built package includes \acs{WasmEdge} support, so crun was compiled from source as part of the \texttt{cloud-init} bootstrap script (guarded by the \texttt{ENABLE\_WASMEDGE=true} flag).
    \item \textbf{containerd plugin entry:} A \texttt{wasmedge} runtime plugin registered in containerd's \texttt{config.toml.tmpl} that invokes crun instead of \texttt{runc}. Pod annotation forwarding (\texttt{run.oci.handler} and \texttt{module.wasm.image/variant}) is configured to allow \acs{K8s} pod annotations to reach crun. This block is appended conditionally when \texttt{ENABLE\_WASMEDGE=true}.
    \item \textbf{WasmEdge~0.14.1:} Installed system-wide at \texttt{/usr/local/lib/libwasmedge.so}. The library is invoked by crun at pod startup to \acs{JIT}-compile the \texttt{.wasm} binary using \acs{LLVM} and execute it.
\end{itemize}

This chain is distinct from the primary SpinKube chain:
\begin{center}
\texttt{kubelet} $\to$ \texttt{containerd} $\to$ \texttt{containerd-shim-spin-v2} $\to$ \textbf{Wasmtime/Cranelift} $\to$ \texttt{.wasm component}
\end{center}

\acs{WasmEdge}~0.14.1 uses an \acs{LLVM} \acs{JIT} backend; Wasmtime/Cranelift (inside Spin) is a different \acs{JIT} backend with different performance characteristics. This backend difference is an acknowledged confounding variable when comparing the P1 and P2 results.

% ── Section 3: WASI P1 Implementation Details ────────────────────────────────
\section{WASI P1 Implementation Details}
\label{app:wasmedge_impl}

\subsection{Wasm/Rust (WasmEdge WASI P1)}

The wasm-rust \acs{WASI}~P1 implementation targeting \texttt{wasm32-wasip1} uses \texttt{wasmedge\_wasi\_socket~0.5.5}~\cite{wasmedgewasisocket}, a \acs{WasmEdge}-specific crate that imports \acs{WasmEdge}'s proprietary \acs{TCP} socket host functions directly from \texttt{wasi\_snapshot\_preview1}. No async \acs{HTTP} framework is available on this target; the server is a single-threaded, synchronous accept-handle-respond loop:

\begin{verbatim}
fn main() {
    let listener = TcpListener::bind("0.0.0.0:8080", false)
        .expect("failed to bind");
    for stream in listener.incoming() {
        match stream {
            Ok(s)  => handle_connection(s),
            Err(e) => eprintln!("accept error: {e}"),
        }
    }
}
\end{verbatim}

Each call to \texttt{handle\_connection} reads the full \acs{HTTP} request, executes the sieve, and writes the response before the loop accepts the next connection. Under concurrent load, incoming connections queue at the \acs{TCP} level.

The \texttt{.wasm} binary is compiled with \texttt{cargo build --target wasm32-wasip1 --release}. The resulting \textbf{OCI image is 0.13\,MB}---16$\times$ smaller than docker-rust.

Source: \texttt{wasm/wasmedge/rust/01-prime-sieve/}

\subsection{Wasm/TinyGo (WasmEdge WASI P1)}

TinyGo~0.40.1 targeting \texttt{wasm32-wasip1} cannot use \texttt{net/http}'s \texttt{ListenAndServe}: TinyGo's Netdev refactor in v0.28 permanently broke the server-side code path. The solution is a custom \texttt{server.go} that directly imports \acs{WasmEdge}'s \acs{TCP} socket host functions via Go's \texttt{//go:wasmimport} directive:

\begin{verbatim}
//go:wasmimport wasi_snapshot_preview1 sock_open
func sock_open(af int32, socktype int32, fd unsafe.Pointer) int32
\end{verbatim}

This calls \acs{WasmEdge}'s \acs{WASI} snapshot preview socket functions (\texttt{sock\_open}, \texttt{sock\_bind}, \texttt{sock\_listen}, \texttt{sock\_accept}, \texttt{sock\_recv}, \texttt{sock\_send}) at the \acs{ABI} level, bypassing the \texttt{net/http} abstraction entirely. The \texttt{main.go} file retains familiar Go \acs{HTTP} handler patterns (\texttt{http.HandleFunc}, \texttt{http.ResponseWriter}), but the transport layer routes through the custom socket implementation.

The \texttt{-gc=conservative} build flag is required: TinyGo's conservative garbage collector scans the stack for pointers without moving objects, ensuring that \texttt{unsafe.Pointer} values passed to the \acs{WasmEdge} socket \acs{ABI} remain valid. The binary is compiled with:

\begin{verbatim}
tinygo build -target=wasi -gc=conservative -opt=2 -o app.wasm .
\end{verbatim}

The resulting \textbf{OCI image is 2.26\,MB}---slightly larger than docker-rust, because TinyGo bundles its own minimal Go runtime (scheduler stub, conservative \acs{GC}, standard library subset) inside the \texttt{.wasm} binary.

Source: \texttt{wasm/wasmedge/tinygo/01-prime-sieve/}

% ── Section 4: WASI P1 Benchmark Results ─────────────────────────────────────
\section{WASI P1 Benchmark Results}
\label{app:wasmedge_results}

The following results were collected on the Hetzner Cloud ccx13 \acs{VM} (2$\times$ AMD EPYC vCPU, 8\,GB RAM) running K3s~v1.35.2+k3s1, WasmEdge~0.14.1, crun~v1.22. The k6 load profile: 1 \acs{VU} ramp to 50~\acp{VU} over 20\,s, hold 60\,s, \texttt{limit=100\,000}.

\subsection{OCI Artifact Size}

\begin{table}[htbp]
    \centering
    \caption{OCI image sizes for the four WASI P1 variants.}
    \label{tab:app_image_sizes}
    \begin{tabular}{ll}
        \hline
        \textbf{Variant} & \textbf{Image Size} \\
        \hline
        wasm-rust     & 0.13\,MB \\
        docker-rust   & 2.08\,MB \\
        wasm-tinygo   & 2.26\,MB \\
        docker-golang & 6.00\,MB \\
        \hline
    \end{tabular}
\end{table}

The wasm-rust image (0.13\,MB) is 16$\times$ smaller than docker-rust and 46$\times$ smaller than docker-golang, reflecting the \acs{Wasm} binary format's elimination of \acs{OS} userland and libc. Unexpectedly, wasm-tinygo (2.26\,MB) is larger than docker-rust because TinyGo bundles its own runtime inside the \texttt{.wasm} binary.

\subsection{Cold-Start and Warm-Start Latency}

\begin{table}[htbp]
    \centering
    \caption{Cold-start latency (N=1; includes image pull) for WASI P1 variants.}
    \label{tab:app_cold_start}
    \begin{tabular}{ll}
        \hline
        \textbf{Variant} & \textbf{Cold Start (ms)} \\
        \hline
        docker-rust   & 38.93 \\
        wasm-tinygo   & 39.98 \\
        wasm-rust     & 42.03 \\
        docker-golang & 3826.40 \\
        \hline
    \end{tabular}
\end{table}

\begin{table}[htbp]
    \centering
    \caption{Warm-start latency over five consecutive runs (image cached). Values in milliseconds.}
    \label{tab:app_warm_start}
    \begin{tabular}{lrrrrr}
        \hline
        \textbf{Variant} & \textbf{Mean} & \textbf{Median} & \textbf{Stdev} & \textbf{Min} & \textbf{Max} \\
        \hline
        docker-rust   &   40.0 &   39.8 &    0.6 &   39.3 &    40.7 \\
        wasm-tinygo   &   40.8 &   41.1 &    1.4 &   38.8 &    42.6 \\
        wasm-rust     &  266.1 &   42.0 &  466.5 &   38.8 &  1098.8 \\
        docker-golang & 5882.1 & 4917.5 & 2323.1 & 3778.6 &  8408.9 \\
        \hline
    \end{tabular}
\end{table}

Three variants cluster tightly between 39\,ms and 42\,ms. docker-golang's 3826\,ms cold start is dominated by its 6\,MB image pull. wasm-rust's warm-start outlier at 1098\,ms (Run~4) reflects a \acs{WasmEdge} \acs{JIT} cache miss: when the \acs{JIT}-compiled native code is not retained between pod restarts, the full \acs{LLVM} compilation pipeline re-executes. docker-golang warm starts (median 4917\,ms) are consistently slow due to Go runtime initialisation under \acs{CPU} contention on the 2-vCPU \acs{VM}.

\subsection{Throughput and Latency Under Load}

\begin{table}[htbp]
    \centering
    \caption{k6 load test results at 50 VUs, 60\,s hold, limit=100\,000.}
    \label{tab:app_k6_results}
    \begin{tabular}{lrrrrr}
        \hline
        \textbf{Variant} & \textbf{RPS} & \textbf{Fail Rate} & \textbf{p50 (ms)} & \textbf{p95 (ms)} & \textbf{p99 (ms)} \\
        \hline
        docker-golang &  509.6 &  0.0\% &   92.4 &   196.3 &   266.3 \\
        docker-rust   &  479.5 &  0.0\% &   92.5 &   157.6 &   173.1 \\
        wasm-rust     &   24.0 & 67.1\% & 5003.0 &  6074.9 &  6266.6 \\
        wasm-tinygo   &    2.9 & 47.9\% & 4605.0 & 18655.8 & 19603.9 \\
        \hline
    \end{tabular}
\end{table}

% ── Section 5: The WASI P1 Concurrency Bottleneck ────────────────────────────
\section{The WASI P1 Concurrency Bottleneck}
\label{app:wasmedge_bottleneck}

The Wasm variants reveal a fundamental architectural constraint. wasm-rust achieves only 24.0\,\acs{RPS} with a 67.1\% failure rate; wasm-tinygo is worse still at 2.9\,\acs{RPS} with 47.9\% failures.

\textbf{Root cause: architectural, not algorithmic.} The \texttt{server\_compute\_us} metric exposes the actual situation. wasm-rust records a median server-side compute time of 117.5\,ms per sieve execution; wasm-tinygo records 243.9\,ms. These are 156$\times$ and 324$\times$ slower than their Docker counterparts for the same algorithm.

Two compounding factors drive this gap:

\begin{enumerate}
    \item \textbf{\acs{WASI}~P1 single-threaded execution.} Both Wasm variants operate a blocking, single-threaded \texttt{accept()}-handle-respond loop. Under 50 concurrent \acp{VU}, incoming connections queue at the \acs{TCP} level. The first accepted request runs to completion before the second is accepted.

    \item \textbf{\acs{WasmEdge} \acs{JIT} overhead for array-intensive workloads.} \acs{Wasm}'s \acl{SFI} model requires that every memory access be bounds-checked at runtime. The sieve algorithm's dense indexed access pattern over 100\,000 elements limits how aggressively the \acs{LLVM} \acs{JIT} can hoist or eliminate bounds checks, producing significant \acs{CPU} overhead per execution compared to native code.
\end{enumerate}

The language-controlled comparison is decisive: wasm-rust (24.0\,\acs{RPS}) vs.\ docker-rust (479.5\,\acs{RPS}) represents a 20$\times$ throughput gap attributable jointly to the \acs{WASI}~P1 single-threaded I/O constraint and \acs{WasmEdge} \acs{JIT} bounds-checking overhead. Disentangling these two factors requires running with \acs{WASI}~P2 multi-threaded I/O---which the primary SpinKube experiment provides.

wasm-tinygo's median server compute time (243.9\,ms) is 2.1$\times$ higher than wasm-rust's (117.5\,ms). TinyGo's conservative \acs{GC} introduces additional overhead: the conservative scan treats every word on the stack as a potential pointer, inflating the cost of the \texttt{[]bool} heap allocation that the sieve performs on each request.

\subsection{Resource Metrics}

\begin{table}[htbp]
    \centering
    \caption{Memory RSS and average CPU utilisation for WASI P1 variants.}
    \label{tab:app_resource_metrics}
    \begin{tabular}{lrrr}
        \hline
        \textbf{Variant} & \textbf{Idle RSS (MB)} & \textbf{Peak RSS (MB)} & \textbf{CPU avg (mcores)} \\
        \hline
        docker-rust   &  2.02 &  2.26 & 299.9 \\
        wasm-rust     &  4.38 &  4.40 & 221.9 \\
        wasm-tinygo   &   --- & 34.92 & 260.3 \\
        docker-golang & 13.63 & 13.63 & 287.4 \\
        \hline
    \end{tabular}
\end{table}

wasm-rust's higher idle \acs{RSS} (4.38\,MB) compared to docker-rust (2.02\,MB) reflects \acs{WasmEdge}'s \acs{JIT} compilation context and linear memory region allocated alongside the process. The wasm-tinygo idle \acs{RSS} was not captured by Prometheus; its peak \acs{RSS} of 34.92\,MB reflects TinyGo's conservative \acs{GC}'s reluctance to compact and release the request-per-allocation boolean array heap.

% ── Section 6: Developer Experience for WasmEdge Variants ────────────────────
\section{Developer Experience Assessment (WasmEdge / WASI P1)}
\label{app:wasmedge_dx}

\begin{table}[htbp]
    \centering
    \caption{Developer experience friction for the WasmEdge WASI P1 variants.}
    \label{tab:app_dx_assessment}
    \adjustbox{max width=\linewidth}{%
    \begin{tabular}{lll}
        \hline
        \textbf{Dimension} & \textbf{wasm-rust (P1)} & \textbf{wasm-tinygo (P1)} \\
        \hline
        Toolchain setup             & Medium    & High      \\
        Dependency compatibility    & High      & Very High \\
        K8s cluster integration     & High      & High      \\
        Debugging and observability & High      & High      \\
        \hline
    \end{tabular}}
\end{table}

\textbf{Dependency compatibility} showed the starkest friction. For wasm-rust, the failure cascade through \texttt{hyper\_wasi}, \texttt{tokio\_wasi}, and \texttt{std::net::TcpListener} required evaluating four candidate \acs{HTTP} approaches before settling on \texttt{wasmedge\_wasi\_socket~0.5.5}~\cite{wasmedgewasisocket}. For wasm-tinygo, writing a custom \texttt{server.go} with \texttt{//go:wasmimport} socket directives and discovering the \texttt{-gc=conservative} requirement required iterative debugging not documented by any official source.

\textbf{Kubernetes integration} required compiling crun from source with \texttt{--with-wasmedge}, a process not available through any managed \acs{K8s} offering (EKS, GKE, AKS). The \texttt{runtimeClassName: wasmedge} and two pod annotations (\texttt{run.oci.handler}, \texttt{module.wasm.image/variant}) are non-standard and not discoverable without reading \acs{WasmEdge}-specific documentation.

\textbf{Debugging} is limited to stdout/stderr output. No debugger can attach to a running \acs{WasmEdge} process; \texttt{kubectl exec} cannot open a shell into a \acs{Wasm} pod. These constraints are eliminated in the primary SpinKube variants, which support the same observability model while providing a more standard deployment path.

% ── Section 7: Optional Deployment ───────────────────────────────────────────
\section{Optional Deployment Instructions}
\label{app:wasmedge_deploy}

The following steps reproduce the WasmEdge/\acs{WASI}~P1 environment alongside the primary SpinKube deployment:

\begin{verbatim}
# 1. In thesis-infra-setup — enable WasmEdge infrastructure:
ENABLE_WASMEDGE=true bash cloud-init.sh
make label-node-wasmedge
make deploy-wasmedge

# 2. Build and push WasmEdge images:
docker build -t docker.io/abdulaziz7225/prime-sieve-wasm-rust:latest \
    wasm/wasmedge/rust/01-prime-sieve/
docker push docker.io/abdulaziz7225/prime-sieve-wasm-rust:latest

docker build -t docker.io/abdulaziz7225/prime-sieve-wasm-tinygo:latest \
    wasm/wasmedge/tinygo/01-prime-sieve/
docker push docker.io/abdulaziz7225/prime-sieve-wasm-tinygo:latest

# 3. Deploy to cluster (optional manifests, NodePorts 30085-30086):
kubectl apply -f k8s/01-prime-sieve/optional/wasm-rust.yaml
kubectl apply -f k8s/01-prime-sieve/optional/wasm-tinygo.yaml
\end{verbatim}

NodePorts: \texttt{wasm-rust} $\to$ 30085, \texttt{wasm-tinygo} $\to$ 30086.
